Large language models confabulate. They produce fluent, confident text that
is factually wrong, and they do so without any detectable hesitation in
their output tokens \citep{ji2023survey}. Current approaches to this problem
fall into two categories: output-level detection (fact-checking the text
after generation) and representation-level probing (training classifiers on
internal activations). The first is slow and requires external knowledge
bases. The second requires model weights and labeled training data, making
it impractical for monitoring deployed systems at scale.

We propose a third approach: \textbf{geometric monitoring of the KV cache
during inference}. The key-value cache is a standard byproduct of
autoregressive generation, used by every transformer backend to avoid
redundant computation. It is already computed, already in memory, and
already accessible through standard inference APIs. We show that singular
value decomposition (SVD) of the key-cache tensor produces features that
distinguish confabulated responses from genuinely uncertain
ones---at inference time, without model weights, without training data,
and with overhead under 50ms per response.

\subsection{A Safety Script, Not a Research Prototype}

The core contribution of this paper is practical: \textbf{a lightweight
confabulation monitor that can be deployed as a drop-in safety layer on
any transformer backend, at any scale.}

The implementation is minimal. After the model generates a response but
before delivering it to the user:

\begin{enumerate}[nosep]
  \item Extract the key cache (already in GPU memory).
  \item Reshape and compute SVD ($< 50$ms on consumer hardware).
  \item Compare spectral features against Marchenko-Pastur thresholds.
  \item Flag, log, or regenerate based on confabulation score.
\end{enumerate}

No fine-tuning. No probing classifier. No access to model internals
beyond the KV cache, which every inference engine (vLLM, TGI,
llama.cpp, TensorRT-LLM) already exposes. The method scales with the
model because the KV cache scales with the model---there is no
additional infrastructure to maintain.

We introduce Marchenko-Pastur (MP) corrected features that further
simplify deployment. Raw SVD features require post-hoc regression to
remove token-count confounds; MP features are approximately
dimension-invariant by construction, requiring no regression, no
calibration data, and no per-model tuning. A single set of thresholds
works across architectures.

\subsection{From Monitoring to Self-Regulation}

Detection alone is a safety tool. Combined with activation steering
and reinforcement learning, it becomes a self-regulation architecture.
The Oracle Loop extends the monitoring approach into a complete cycle:

\begin{enumerate}[nosep]
  \item \textbf{Detect}: Extract spectral features from the generation-phase
    KV cache using MP-corrected SVD.
  \item \textbf{Decide}: Compare geometry against encoding-phase baseline to
    classify cognitive state.
  \item \textbf{Steer}: On misalignment detection, roll back to the
    encoding-phase cache checkpoint and regenerate with activation
    addition.
  \item \textbf{Train}: Feed geometry-derived rewards into GRPO training,
    teaching the model to avoid confabulation geometry \emph{from the
    inside}.
\end{enumerate}

Steps 1--2 are the safety monitor. Step 3 adds runtime correction.
Step 4 is the training loop that makes the runtime monitor progressively
less necessary. This paper focuses primarily on the monitor (Steps 1--2)
and describes the full architecture (Steps 3--4) as future-facing work
that is architecturally complete but not yet converged.

\subsection{Contributions}

This paper makes four contributions:

\begin{enumerate}[nosep]
  \item \textbf{A deployable confabulation monitor} based on KV-cache SVD
    that works on any transformer backend without model weights or
    training data (\cref{sec:methods}).
  \item \textbf{Marchenko-Pastur corrected features} that eliminate
    the token-count confound through random matrix theory rather than
    post-hoc regression (\cref{sec:methods}).
  \item \textbf{A clean detection experiment} with post-hoc behavioral
    classification, generation-only cache analysis, and a 14-control
    battery, across three architectures (\cref{sec:results}).
  \item \textbf{The Oracle Loop architecture}: a six-layer harness for
    real-time inference monitoring with geometry-aware GRPO training
    (\cref{sec:architecture}).
\end{enumerate}

\subsection{Scope and Honest Framing}

The detection results are moderate (AUROC 0.707 on the hardest
comparison), not perfect. This is a confabulation \emph{signal}, not a
confabulation \emph{oracle}. But it is a signal that requires no training,
no labeled data, no weight access, and under 50ms of compute---properties
that no competing approach offers simultaneously. We believe the
combination of accessibility, generality, and low overhead makes this
useful even at moderate accuracy, particularly as one signal among several
in a defense-in-depth safety stack.
